\documentclass[12pt]{article}
\usepackage[a4paper,margin=1in]{geometry}
\usepackage{amsmath,amssymb}
\usepackage{mathtools}
\usepackage{bm}

\title{Numerically Stable Coherent-State Hilbert-Space Truncation}
\author{}
\date{}

\begin{document}
\maketitle

\section*{Purpose}
For the single-site coherent-state runs in this project, we choose the smallest Hilbert-space truncation that preserves the coherent-state mean particle number to within a prescribed relative tolerance. The original implementation used a direct finite sum containing factorials. That formula is physically correct, but it becomes numerically unstable for large mean occupation. This note records the original formula, derives a better equivalent expression, and explains why the new algorithm is preferred.

\section*{Original Truncation Criterion}
Let
\[
\mu = |\alpha|^2
\]
denote the coherent-state mean particle number. The coherent-state number distribution is Poisson:
\[
P(n) = e^{-\mu}\frac{\mu^n}{n!}, \qquad n=0,1,2,\dots
\]

If the occupation basis is truncated at a maximum occupation $x$, then the truncated mean particle number is
\[
\langle n \rangle_x
= e^{-\mu}\sum_{n=0}^{x} n\,\frac{\mu^n}{n!}.
\]

In the project, the full coherent-state mean is
\[
\langle n \rangle_\infty = \mu,
\]
and the truncation rule is:
\[
\left|\langle n \rangle_x - \mu\right| \le \varepsilon \mu,
\]
where $\varepsilon$ is the adjustable relative tolerance, currently $0.01$ for a $1\%$ criterion.

The Hilbert-space size is then
\[
d = x+1,
\]
because the basis includes occupations $0,1,\dots,x$.

\section*{Why the Original Direct Sum Is Problematic}
The direct implementation evaluates terms of the form
\[
e^{-\mu}\frac{\mu^n}{n!}.
\]
For large $\mu$ and large $n$, this is a poor numerical strategy because:
\begin{itemize}
\item $\mu^n$ becomes enormous and can overflow.
\item $n!$ grows extremely rapidly and can also overflow.
\item Even when both quantities are representable individually, their ratio may suffer severe loss of numerical accuracy.
\item The implementation has to repeat this unstable computation many times while searching for the smallest acceptable cutoff.
\end{itemize}

In practice, this is exactly what breaks for large coherent occupations such as $\mu = 1000$ and beyond.

\section*{Key Identity for the Truncated Mean}
We now derive an equivalent formula that avoids direct factorial summation in the mean calculation.

Start from
\[
\langle n \rangle_x = \sum_{n=0}^{x} nP(n).
\]
Since the $n=0$ term vanishes, we write
\[
\langle n \rangle_x = \sum_{n=1}^{x} n\,e^{-\mu}\frac{\mu^n}{n!}.
\]
Now use
\[
n\frac{\mu^n}{n!} = \mu\frac{\mu^{n-1}}{(n-1)!}.
\]
Therefore
\[
\langle n \rangle_x
= \mu e^{-\mu}\sum_{n=1}^{x}\frac{\mu^{n-1}}{(n-1)!}.
\]
Let $m=n-1$. Then
\[
\langle n \rangle_x
= \mu e^{-\mu}\sum_{m=0}^{x-1}\frac{\mu^m}{m!}.
\]
Recognizing the Poisson cumulative distribution function,
\[
F_{\text{Pois}}(k;\mu)
= \sum_{m=0}^{k} e^{-\mu}\frac{\mu^m}{m!},
\]
we obtain the exact identity
\[
\boxed{
\langle n \rangle_x = \mu\,F_{\text{Pois}}(x-1;\mu)
}
\qquad \text{for } x \ge 1.
\]

\section*{Tail-Error Formulation}
The truncation error in the mean is
\[
\mu - \langle n \rangle_x
= \mu\left(1 - F_{\text{Pois}}(x-1;\mu)\right).
\]
Using the Poisson survival function
\[
S_{\text{Pois}}(k;\mu) = 1 - F_{\text{Pois}}(k;\mu),
\]
we get
\[
\boxed{
\mu - \langle n \rangle_x
= \mu\,S_{\text{Pois}}(x-1;\mu)
}.
\]

Hence the original relative-tolerance rule
\[
\left|\langle n \rangle_x - \mu\right| \le \varepsilon\mu
\]
is equivalent to
\[
\mu\,S_{\text{Pois}}(x-1;\mu) \le \varepsilon\mu,
\]
and for $\mu>0$ this simplifies to
\[
\boxed{
S_{\text{Pois}}(x-1;\mu) \le \varepsilon
}.
\]

This is the most important result: instead of summing the truncated first moment directly, we only need the Poisson tail probability.

\section*{Why the New Formulation Is Better}
The new formulation is preferable for several reasons.

\subsection*{1. It is mathematically equivalent}
We did not change the physics criterion. We only rewrote the same truncated-mean condition in a more stable form. The selected Hilbert-space size still means:
\begin{quote}
Choose the smallest occupation cutoff whose missing contribution to the coherent-state mean is below the chosen relative tolerance.
\end{quote}

\subsection*{2. It avoids direct factorial overflow}
The old implementation explicitly evaluated $\mu^n/n!$ term by term. The new implementation calls a numerically stable Poisson CDF / survival-function backend from SciPy, which is specifically designed for this type of probability calculation.

\subsection*{3. It gives a monotone search criterion}
As $x$ increases, the Poisson tail $S_{\text{Pois}}(x-1;\mu)$ decreases monotonically. Therefore the acceptance criterion is monotone:
\[
x \uparrow \quad \Longrightarrow \quad S_{\text{Pois}}(x-1;\mu)\downarrow.
\]
This means we can use a reliable search strategy:
\begin{enumerate}
\item Start from a reasonable upper guess.
\item Double it until the tolerance is satisfied.
\item Use binary search to find the smallest acceptable cutoff.
\end{enumerate}
This is much cleaner than incrementing one occupation at a time with an unstable inner sum.

\subsection*{4. It scales to much larger coherent occupations}
The old implementation breaks for large $\mu$ because of overflow in the direct Poisson-series terms. The new implementation is stable for large coherent-state occupations because it relies on special-function evaluations rather than raw factorial arithmetic.

\section*{Final Algorithm Used in Code}
Let $\mu = \texttt{num\_of\_particles}$ and let $\varepsilon$ be the relative tolerance.

\subsection*{Step 1: Define the acceptance test}
For a candidate maximum occupation $x$, compute the missing mean
\[
\Delta_x = \mu\,S_{\text{Pois}}(x-1;\mu).
\]
Accept the truncation when
\[
\Delta_x \le \varepsilon \mu.
\]

\subsection*{Step 2: Find an upper bound}
Start with
\[
x_{\text{high}} = \max(1,\lceil \mu \rceil).
\]
If this is not sufficient, repeatedly double it until the criterion is satisfied.

\subsection*{Step 3: Binary search}
Once we have:
\begin{itemize}
\item a lower cutoff that fails, and
\item an upper cutoff that passes,
\end{itemize}
apply binary search to obtain the smallest acceptable $x$.

\subsection*{Step 4: Convert to Hilbert-space size}
The Hilbert-space size is
\[
d = x+1.
\]

\section*{Why This Is the Best Choice Here}
Among the possible alternatives, this approach is the best fit for the project.

\subsection*{Compared with direct factorial summation}
It is vastly more stable and avoids overflow.

\subsection*{Compared with a recursive Poisson-term summation}
Recursive summation is better than raw factorials, but it is still less direct than using the exact CDF / survival-function identity for the truncation criterion.

\subsection*{Compared with a normal approximation}
A normal approximation could provide a rough estimate for large $\mu$, but it changes the criterion from exact Poisson-tail control to an approximation. Since SciPy is already available, there is no reason to give up the exact Poisson-based criterion here.

\section*{Conclusion}
The project now uses a numerically stable coherent-state truncation algorithm based on the Poisson survival function. The new method:
\begin{itemize}
\item preserves the original physical tolerance rule,
\item avoids factorial overflow,
\item supports large coherent occupations much more reliably,
\item and allows an efficient monotone search for the smallest acceptable Hilbert-space size.
\end{itemize}

This makes it the most appropriate formulation for future larger-scale coherent-state runs in the project.

\end{document}
